\chapter{Fazit und Ausblick}\label{ch:conclusion}

Dieses Kapitel zieht Bilanz: Es stellt je Forschungsfrage FF0--FF4 (Abschnitt~\ref{sec:rqs}) die
Teilfragen voran, benennt den Stand ihrer Beantwortung in drei Stufen --- beantwortet, strukturell
beantwortet, offen ---, führt je Frage die Beweiskette durch die Kapitel und ihre Folgen, ordnet die
Hypothesen dem Zielpfad zu, benennt die Grenzen des Systems aus Sicht des Anwenders und skizziert die
nächsten Schritte.

\section{Beantwortung der Forschungsfragen}\label{sec:answers}
Die Antworten folgen dem Schema Behauptung, Beleg, Validierung~\cite{shaw2003writing}: Eine Frage gilt
als \emph{beantwortet}, wenn ein Test oder eine Messung sie belegt; als \emph{strukturell
beantwortet}, wenn der Apparat, der die Antwort liefert, implementiert und getestet ist, die Messung aber
nicht vorliegt; als \emph{offen}, wenn weder Struktur noch Messung vorliegen. Drei Belegarten kommen vor: die
Struktur (ein Code-Gegenstand, der die Antwort stützt), der Test (eine statische Prüfung oder ein
Vertrags-Test, der sie sichert) und die Messung (eine Zahl mit ihrem Lauf); für die Frage nach dem
Stand der Technik der Workloads (FF2.d) tritt der Literaturbeleg hinzu. Ein \enquote{strukturell
beantwortet} ersetzt die Zahl nicht: Ohne Messwerte bleibt die Hauptfrage unbeantwortet, auch wenn ihre
Struktur tragfähig ist. Tabelle~\ref{tab:ff-beweiskette} führt Status, Beweis und Rest je Haupt- und Teilfrage;
die Berichtsform folgt den Regeln der empirischen Software- und
Systemforschung~\cite{wohlin2012experimentation,runeson2009guidelines} und der Forderung, Bezugsgröße,
Plattform und nicht erhobene Werte an jedem Ergebnis auszuweisen~\cite{blackburn2016truth,hoefler2015benchmarking}.

\begin{nurlang}
%% INTERN-272: KA-112
Quer zu allen fünf Fragen hat die Einleitung eine Kategorisierungs-Frage als Frage ausgelegt: ob sich die Achsen
des Entwurfsraums sauber in Organ-, System- und Mess-Achsen trennen lassen --- und wie eine solche
Trennung zu bewerkstelligen ist, wo sie in einem Entwurf nicht angelegt ist. Beides ist beantwortet
(Abschnitt~\ref{sec:axis-triad}): Die Achsen trennen sich entlang der drei Belange in drei \emph{Realms}
--- Organ, System, Mess --- mit je eigener Klassen-Wurzel, je eigenem Angebots-Katalog (drei generierte
Registries: die Organ-Achsen-Registry der Cache-Engine, die System-Achsen-Registry und die
Mess-Achsen-Registry) und je eigener Stempel-Zeile an der Binary (Abschnitt~\ref{sec:stamp-model}); die
drei Zeilen-Träger des Fingerprints --- Organ-, System- und Mess-Zeile --- sind als Typen paarweise
unkonstruierbar auseinander, sodass eine Vertauschung nicht übersetzt. In die \texttt{binary\_id} gehen
ausschließlich die Organ-Haupt-Achsen ein; System- und Mess-Achsen sind in ihren Registries als
binary-id-neutral deklariert. Gehalten wird die Trennung durch übersetzungszeitliche
\emph{Rückkehr-Sperren} (Abschnitt~\ref{sec:impl-system-axes}): \texttt{static\_assert}-Prüfungen weisen
Compiler, Scheduling und Last-Rahmenwerk als System-Haupt-Achsen zurück --- das Scheduling ist eine
Unter-Achse der Ziel-ISA \texttt{target\_isa} mit fünf Dimensionen, der Compiler eine Unter-Achsen-Gruppe
der Komplex-Klammer \texttt{build\_target\_complex}, das Last-Rahmenwerk die Meta-Meta-Achse des
Mess-Realms ---, und die Flag-Grammatik der Versions-Stempel lässt nur katalogisierte Flag-Ketten zu.
Diese Dreiteilung ist die zweite Gestalt des Trennbarkeits-Problems und stützt die Antworten auf die
folgenden Teilfragen.
\end{nurlang}
\begin{nurkurz}
%% INTERN-272: KA-112
Quer zu allen fünf Fragen hat die Einleitung die Kategorisierungs-Frage ausgelegt, ob sich die Achsen des
Entwurfsraums sauber in Organ-, System- und Mess-Achsen trennen lassen --- und wie eine solche Trennung zu
bewerkstelligen ist, wo sie in einem Entwurf nicht angelegt ist. Beides ist beantwortet
(Abschnitt~\ref{sec:axis-triad}): Die Achsen trennen sich in drei \emph{Realms} --- Organ, System, Mess ---
mit je eigener Klassen-Wurzel, generierter Registry und Stempel-Zeile an der Binary
(Abschnitt~\ref{sec:stamp-model}); die drei Zeilen-Träger des Fingerprints sind als Typen paarweise
unkonstruierbar auseinander. In die \texttt{binary\_id} gehen nur die Organ-Haupt-Achsen ein, System- und
Mess-Achsen sind binary-id-neutral. Übersetzungszeitliche \emph{Rückkehr-Sperren}
(Abschnitt~\ref{sec:impl-system-axes}) weisen Compiler, Scheduling und Last-Rahmenwerk als
System-Haupt-Achsen
zurück: Das Scheduling ist Unter-Achse der Ziel-ISA, der Compiler Unter-Achsen-Gruppe der Komplex-Klammer,
das Last-Rahmenwerk die Meta-Meta-Achse des Mess-Realms. Diese Dreiteilung ist die zweite Gestalt des
Trennbarkeits-Problems und stützt die Antworten auf die folgenden Teilfragen.
\end{nurkurz}

\begin{nurlang}
%% INTERN-272: KA-113
\paragraph{FF0 (Hauptfrage).} Gefragt ist, wie stark aktive, messgetriebene cache-bewusste
Entscheidungen --- Seitentyp, Speicher-Layout, Value-Ablage, Prefetching --- ihre passiven, statisch
ausgelegten Gegenstücke auf unterschiedlichen CPU-Plattformen übertreffen; die Erweiterung FF0.E fragt,
ob sich ein gemessener Gewinn eindeutig dem Organ-, dem System- oder dem Mess-Belang zuordnen lässt.
Der \emph{Status}: FF0 ist offen. Ohne die End-to-End-Läufe je Plattform gibt es keine Zahl; die Struktur,
die die Zahl liefert, ist implementiert. Dieselbe achsen-konstante Struktur wird einmal aktiv --- mit
laufzeit-dynamischer Cache-Line-Politik (der Politik-Selektor leitet aus einem ausgewerteten
Profil-Aggregat aus Operationsmix und Working-Set die Cache-Line-Anpassung ab) und messwert-getriebener
Auswahl der besten Binary --- und einmal passiv --- mit statisch fixierter Organ-Wahl von
Speicher-Layout, Prefetch, Value-Ablage und Knotentyp als Segmenten der \texttt{binary\_id} --- gebaut
und unter identischer Compiler- und Flag-Basis gegenübergestellt (Standard \texttt{-O2};
\texttt{-O3} nur als Opt-in unter Warnung und mit eigener Toolchain-Identität). Vermessene Messmaschinen sind
die drei registrierten Maschinen-Signaturen: prod1 (AMD Zen~5, Ryzen~9~9950X3D), prod2 (Intel Core % chktex 29
i9-12900K, Alder Lake --- eine Hybrid-CPU mit P- und E-Kernen, ohne AVX-512) und ein ODROID-Board mit
Intel-Gracemont-Kernen; die in der Frage genannte Server-Plattform Sapphire Rapids ist in der
Plattform-Erkennung angelegt (ZIH-Knoten Barnard, Xeon~8470), aber nicht gemessen (n/a;
Abschnitt~\ref{sec:eval-platforms}). FF0.E ist strukturell beantwortet, und die dafür nötige Trennung ist
die Dreiteilung: Ein Organ-Effekt permutiert die \texttt{binary\_id}, ein Plattform-Belang steht als
System-Stempel-Zeile und Bau-Versionskennung an der Binary, ohne sie zu permutieren, und ein Mess-Belang
bleibt binary-id-neutral --- kein Mess-Instrument verändert, \emph{was} gebaut wird. Als
\emph{Beleg} dienen Abschnitt~\ref{sec:axis-triad} (Dreiteilung; Struktur),
Abschnitt~\ref{sec:impl-system-axes}
(Rückkehr-Sperren; Test) und Abschnitt~\ref{sec:eval-method} (Messplan; Struktur). Die \emph{Folge}: Eine
Bezifferung setzt den vollständigen Bau und den Lauf der Reihe~A voraus (Abschnitt~\ref{sec:outlook}); der
Effekt der Cache-Line-Bewusstheit aus den synthetisierten Mess-Kurven ist der erste Auswertungsschritt
nach dem ersten Messlauf (Abschnitt~\ref{sec:heuristics-curves}).
\end{nurlang}
\begin{nurkurz}
%% INTERN-272: KA-113
\paragraph{FF0 (Hauptfrage).} Gefragt ist, wie stark aktive, messgetriebene cache-bewusste Entscheidungen
ihre passiven, statisch ausgelegten Gegenstücke auf unterschiedlichen CPU-Plattformen übertreffen und ob
sich
ein gemessener Gewinn dem Organ-, System- oder Mess-Belang zuordnen lässt (FF0.E). Der \emph{Status}: FF0
ist
offen --- ohne End-to-End-Läufe keine Zahl; die Struktur ist implementiert: Dieselbe Struktur wird einmal
aktiv (laufzeit-dynamische Cache-Line-Politik, die der Politik-Selektor aus einem ausgewerteten
Profil-Aggregat aus Operationsmix und Working-Set ableitet, und
messwert-getriebene Binary-Auswahl) und einmal passiv (statisch fixierte Organ-Wahl in der
\texttt{binary\_id}) gebaut und unter identischer Compiler- und Flag-Basis verglichen (Standard
\texttt{-O2}, \texttt{-O3} nur als Opt-in). Messmaschinen sind prod1, prod2 und ein ODROID-Board; Sapphire
Rapids ist angelegt, nicht gemessen (n/a; Abschnitt~\ref{sec:eval-platforms}). FF0.E ist strukturell
beantwortet durch die Dreiteilung: Ein Organ-Effekt permutiert die \texttt{binary\_id}, ein
Plattform-Belang steht als System-Stempel-Zeile an der Binary, ein Mess-Belang bleibt binary-id-neutral.
\emph{Beleg}: Abschnitt~\ref{sec:axis-triad} (Struktur), Abschnitt~\ref{sec:impl-system-axes} (Test),
Abschnitt~\ref{sec:eval-method} (Messplan).
\emph{Folge}: Die Bezifferung setzt den vollständigen Bau und den Lauf der Reihe~A voraus
(Abschnitt~\ref{sec:outlook}); der Cache-Line-Effekt aus den synthetisierten Mess-Kurven folgt als erster
Auswertungsschritt (Abschnitt~\ref{sec:heuristics-curves}).
\end{nurkurz}

\begin{nurlang}
%% INTERN-272: KA-115
\paragraph{FF1 (Komposition).} Gefragt ist, wie sich ein Suchalgorithmus-Entwurf kanonisch in
orthogonale Achsen zerlegen und in einer Permutationsmatrix systematisieren lässt, sodass
Stand-der-Technik-Verfahren möglichst aus ihrem Originalcode bias-frei als Profile rekonstruierbar
und beliebig permutierbar werden; FF1.a
fragt nach dem Kategorien-Schnitt der Zerlegung, FF1.b nach den Organ-Kopplungen und dem Verhalten unter
allein variierter ISA-/SIMD-Achse. Der \emph{Status}: FF1 ist strukturell beantwortet, FF1.a beantwortet,
FF1.b in seinem empirischen Teil offen. Die kanonische Zerlegung ist tragfähig: \emph{Achtzehn}
Organ-Haupt-Achsen in fester Slot-Reihenfolge T0--T17 --- als konstantes Namens-Array der
Achsen-Pfad-Serialisierung Code, nicht Konvention --- spannen die Permutationsmatrix auf
(Abschnitt~\ref{sec:axis-triad}); dazu tritt eine optionale \emph{Organ-Meta-Meta-Achse} für die
Festplatten-Ein-/Ausgabe \texttt{disk\_io}, in der Organ-Registry als neunzehnter Eintrag geführt
(übersetzungszeitliche Stufe mit der Laufzeit-Option \texttt{staging}, binary-id-neutral), per
Experiment-Profil aber nicht wählbar, weil der Profil-Parser ihren Registry-Block nicht liest; der
Profil-Anschluss ist vorgesehen, aber nicht implementiert. Sie ist das Beispiel der Arbeit für die
Meta-Meta-Ebene M3, die Indirektion einer Beschreibung: Sie stempelt im Klammer-Anhang der Organ-Zeile,
und die Klammer-Tiefe eines Eintrags ist seine Meta-Meta-Ebene (Abschnitt~\ref{sec:stamp-model}). Zwei
Bezugsgrößen sind zu trennen: Der golden-Referenzkatalog variiert siebzehn Achsen je zwei und legt das
Persistenz-Ziel auf einen Wert fest --- $2^{17} = 131072$ Binary-Identitäten je System-Permutation als
Regressions-Detektor mit einer festgeschriebenen CRC-64-Prüfsumme\footnote{CRC-64: eine zyklische
Redundanzprüfung mit 64-Bit-Prüfsumme (hier die Variante ECMA-182), berechnet über alle Identitäten in
kanonischer Reihenfolge; eine geänderte Identität ändert die Prüfsumme.} über alle Identitäten
(Abschnitt~\ref{sec:impl-contracts}) ---, während die Mächtigkeit eines geplanten Raums keine Konstante
ist, sondern die Rechnung des Planers je Profil und Plattform (Subkommando \texttt{simulate}); alle
achtzehn Achsen sind für die Variation vorgesehen (nicht gemessen). Die 33 Korpus-Arbeiten P01--P33 besitzen je ein
Profil-XML, davon 30 Vollprofile eigenständiger Suchentwürfe (acht Rang-1 und 22 Rang-2/3) und drei
abstrakte Profile (OLC, LOUDS, VAMPIR), die nur einzelne Achsen belegen
(Abschnitt~\ref{sec:sota-instances}); die Registry der Paper-Prüflinge hält diese Übersetzung
übersetzungszeitlich gegen Dubletten und Lücken. Die variadische Abbildung --- ein Template-Parameter für
die \texttt{std::vector}-API, zwei für \texttt{std::map}, $N>2$ für
\texttt{std::map<K,\,std::tuple<\ldots>>} (Anhang~\ref{app:interfaces}) --- macht algorithmus- wie
container-orientierte Muster auf einer Engine ausdrückbar, und die Gattungs-Dreiteilung Map, Container,
Graph ordnet die Genera darunter (Abschnitt~\ref{sec:anatomy-levels}). FF1.a ist bejaht: Die Zerlegung
bleibt nur orthogonal, weil Plattform- und Mess-Belange aus der Permutationsmatrix herausgehalten sind,
und die Rückkehr-Sperren halten diesen Schnitt (Abschnitt~\ref{sec:impl-system-axes}). Zu FF1.b: Die von
der Frage genannte Kopplung Knotentyp $\times$ SIMD $\times$ Pfadkompression zerfällt --- Knotentyp
$\times$ Pfadkompression bleibt eine nicht voll-orthogonale Organ-Kopplung, die als solche markiert und
bei der Permutation eingeschränkt wird, während SIMD keine Organ-Achse ist, sondern eine Bau-Steuergröße
der System-Achse \texttt{external\_utils} (Laufzeit-Unter-Achse mit den Optionen
\texttt{no\_extension}, \texttt{avx2} und \texttt{avx512} und einem Merkmals-Katalog von 23 Flags,
binary-id-neutral): Ein und dieselbe \texttt{binary\_id} wird unter mehreren System-Permutationen
gebaut und bei sonst unveränderter Struktur verglichen; die empirische Antwort zum Cache-Update- und
Cache-Line-Füllverhalten ist eine Messung und liegt nicht vor. Als \emph{Beleg} dienen
Abschnitt~\ref{sec:axis-triad} (Zerlegung und Realms; Struktur), Abschnitt~\ref{sec:mess-chain}
(Katalog
und Bezugsgröße; Struktur) und Abschnitt~\ref{sec:impl-contracts} (Katalog-Prüfungen und Golden-Prüfsumme; Test).
Die \emph{Folge}: FF1 strukturell beantwortet; FF1.b (Füllverhalten) setzt den vollständigen Bau voraus.
\end{nurlang}
\begin{nurkurz}
%% INTERN-272: KA-115
\paragraph{FF1 (Komposition).} Gefragt ist, wie sich ein Suchalgorithmus-Entwurf kanonisch in orthogonale
Achsen zerlegen und in einer Permutationsmatrix systematisieren lässt, sodass Stand-der-Technik-Verfahren
bias-frei als Profile rekonstruierbar und permutierbar werden; FF1.a fragt nach dem Kategorien-Schnitt,
FF1.b
nach den Organ-Kopplungen und der allein variierten ISA-/SIMD-Achse. Der \emph{Status}: FF1 ist strukturell
beantwortet, FF1.a beantwortet, FF1.b empirisch offen.
\emph{Achtzehn} Organ-Haupt-Achsen in fester Slot-Reihenfolge T0--T17 (Namens-Array im Code) spannen die
Matrix auf (Abschnitt~\ref{sec:axis-triad}); \texttt{disk\_io} ist neunzehnter Registry-Eintrag und Beispiel
der Meta-Meta-Ebene M3: gestempelt im Klammer-Anhang der Organ-Zeile, die Klammer-Tiefe ist die
Meta-Meta-Ebene (Abschnitt~\ref{sec:stamp-model}). Zwei Bezugsgrößen sind zu trennen: der
golden-Referenzkatalog mit $2^{17} = 131072$ Binary-Identitäten je System-Permutation, per CRC-64-Prüfsumme
gegen Regression gesichert (Abschnitt~\ref{sec:impl-contracts}), und die vom Planer gerechnete Mächtigkeit
eines geplanten Raums (\texttt{simulate}). Die 33 Korpus-Arbeiten besitzen je ein Profil-XML (30
Vollprofile,
drei abstrakte; Abschnitt~\ref{sec:sota-instances}); die variadische Abbildung (Anhang~\ref{app:interfaces})
trägt algorithmus- wie container-orientierte Muster, die Gattungs-Dreiteilung Map, Container, Graph ordnet
die Genera darunter (Abschnitt~\ref{sec:anatomy-levels}).
FF1.a ist bejaht: Die Zerlegung bleibt orthogonal, weil Plattform- und Mess-Belange herausgehalten sind, und
die Rückkehr-Sperren halten den Schnitt (Abschnitt~\ref{sec:impl-system-axes}). Zu FF1.b: Die Kopplung
Knotentyp $\times$ SIMD $\times$ Pfadkompression zerfällt --- Knotentyp $\times$ Pfadkompression bleibt eine
markierte Organ-Kopplung, SIMD ist Bau-Steuergröße der System-Achse \texttt{external\_utils}; die Messung
des
Füllverhaltens liegt nicht vor.
\emph{Beleg}: Abschnitte~\ref{sec:axis-triad} und~\ref{sec:mess-chain} (Struktur),
Abschnitt~\ref{sec:impl-contracts} (Test).
\emph{Folge}: FF1 strukturell beantwortet; FF1.b setzt den vollständigen Bau voraus.
\end{nurkurz}

\begin{nurlang}
%% INTERN-272: KA-116
\paragraph{FF2 (Mess-Methodik).} Gefragt ist, wie eine reproduzierbare Mess-Pipeline mit
Hardware-Zähler-Erhebung gestaltet sein muss, damit jede Achsen-Permutation fair --- unter gleicher
Compiler- und Flag-Basis --- in drei Messreihen über drei Granularitäten gemessen wird, und unter welchen
plattform-kalibrierten Kriterien eine Plattform als verifiziert gilt; FF2.a fragt, ob die
Mess-Instrumente selbst eigene Achsen sein müssen, FF2.b nach der Nachmessung widersprüchlicher
Literaturbefunde, FF2.c nach der Erweiterbarkeit, FF2.d nach den Workloads, die gerade für
Suchalgorithmen zum Stand der Technik gehören. Der \emph{Status}: FF2 ist strukturell beantwortet,
FF2.a und FF2.d beantwortet. Die Messung ist in
einen Produktiv- und einen Experiment-Modus geteilt, geschaltet über den Master-Schalter der
Mess-Übersetzung (\texttt{COMDARE\_\allowbreak{}MEASUREMENT\_\allowbreak{}ON}): Der Experiment-Modus
misst stets --- ein Release-Bau ohne Messfühler ist ein eigenes, auslieferbares Artefakt, das den
Schalter erzwungen abschaltet, kein abgeschalteter Messlauf ---, und im Lager wird nur nachgemessen, was
dort fehlt; das übersetzungsstatische Schlüssel-Schema der vier Bestände ist der Entwurf dieser
Wiedererkennung, der Lager-Vollausbau ist nicht implementiert (Abschnitt~\ref{sec:lager}). Den
Ablauf steuert der \texttt{ExperimentDriver} host-seitig: Der Runner lädt die Module und führt die Phasen
Laden, Ausführen, Messen und Abbau, der Treiber führt Ausführung und Messung je Modul und sammelt die
Ergebnisse im \texttt{ResultAggregator}. Jede Reihe wird auf drei \emph{Mess-Ebenen} erhoben, in genau
dieser Reihenfolge: der Wallclock-Ebene (\texttt{wallclock}; die Gesamtläufe im CacheEngineBuilder über
die Last-Rahmenwerke), der Makro-Ebene (\texttt{macro}; die Gattungs- und Genus-Interface-Operationen
der Tier-Binary) und der Mikro-Ebene (\texttt{micro}; Achsen-Interface und Erfolgs-Parameter); w/ma/mi
ist der registrierte Alias, und die drei Ebenen sind die Mess-Tooling-Haupt-Achse, die je Wahl in eine
CacheEngineBuilder-Strecke einkompiliert wird (Abschnitt~\ref{sec:mess-chain}). FF2.a ist bejaht: Die
Mess-Instrumente sind ein eigener Realm im Planer --- das Last-Rahmenwerk als Meta-Meta-Achse, das
Mess-Tooling als Haupt-Achse, die Mess-Kategorien als Unter-Achse (sechzehn Bausteine in der Registry;
vorgesehen ist die Erhebung je Achsen-Interface-Parametersatz, deren Katalog je Achse nicht
implementiert ist) ---, stets präsent, sobald gemessen wird, und binary-id-neutral. Bewerkstelligt wird
die Trennung durch den doppelten Wurzel-Schnitt: Der System-Realm wurzelt am CacheEngineBuilder, der
Mess-Realm am Planer, und die Mess-Achse treibt über die System-Achsen-Wurzel die Freigabe der
Mess-Quelle in Gestalt des CacheEngineBuilder-Typs --- die Legenden-Kette Mess-Achse $\to$
CacheEngineBuilder-Typ $\to$ System-Achsen $\to$ Tier-Binary ist als Namensschema Code
(Abschnitt~\ref{sec:axis-triad}). Beim Workload reicht der Treiber einen Standard-Workload je Modul
durch, und ein Profil bringt seinen eigenen mit; vorgesehen ist die Voll-Permutation über alle
verfügbaren Workloads, damit kein Entwurf nur unter seinem selbstgewählten Lastprofil gewinnt --- der
\emph{Lastprofil-Selektionsverzerrung}\footnote{Lastprofil-Selektionsverzerrung (workload selection
bias): die Verzerrung eines Vergleichs dadurch, dass die Lastprofile so gewählt sind, dass ein Entwurf
unter ihnen bevorzugt gut abschneidet; die Benchmark-Literatur verlangt darum eine breite, begründete
Lastprofil-Menge~\cite{blackburn2006dacapo,blackburn2016truth}.} begegnet der Apparat mit der Vollmatrix
aller Lastprofile gegen alle Lebewesen-Binaries, weil schon unauffällige Setup-Entscheidungen
Messergebnisse systematisch verfälschen~\cite{mytkowicz2009wrongdata}; Umsetzungsstand: Die
Voll-Permutation ist nicht implementiert.
\end{nurlang}
\begin{nurkurz}
%% INTERN-272: KA-116
\paragraph{FF2 (Mess-Methodik).} Gefragt ist, wie eine reproduzierbare Mess-Pipeline mit
Hardware-Zähler-Erhebung jede Achsen-Permutation fair in drei Messreihen über drei Granularitäten misst und
wann eine Plattform als verifiziert gilt; FF2.a fragt nach Mess-Instrumenten als eigenen Achsen, FF2.b nach
widersprüchlichen Literaturbefunden, FF2.c nach der Erweiterbarkeit, FF2.d nach den Workloads des Stands der
Technik. Der \emph{Status}: FF2 ist strukturell beantwortet, FF2.a und FF2.d beantwortet. Produktiv- und
Experiment-Modus schaltet der Master-Schalter \texttt{COMDARE\_\allowbreak{}MEASUREMENT\_\allowbreak{}ON};
der
Experiment-Modus misst stets, im Lager wird nur Fehlendes nachgemessen (Vollausbau nicht implementiert;
Abschnitt~\ref{sec:lager}); der \texttt{ExperimentDriver} steuert die Runner-Phasen Laden, Ausführen,
Messen,
Abbau. Jede Reihe wird auf drei \emph{Mess-Ebenen} erhoben --- \texttt{wallclock}, \texttt{macro},
\texttt{micro} ---; w/ma/mi ist der registrierte Alias, die drei Ebenen sind die Mess-Tooling-Haupt-Achse
(Abschnitt~\ref{sec:mess-chain}). FF2.a ist bejaht:
Die Mess-Instrumente sind ein eigener Realm im Planer --- das Last-Rahmenwerk \texttt{load\_framework} als
Meta-Meta-Achse, das Mess-Tooling als Haupt-Achse, die sechzehn Mess-Kategorien als Unter-Achse ---,
stets präsent und binary-id-neutral; die Trennung hält der doppelte Wurzel-Schnitt
(System-Realm am CacheEngineBuilder, Mess-Realm am Planer), und die Legenden-Kette Mess-Achse $\to$
CacheEngineBuilder-Typ $\to$ System-Achsen $\to$ Tier-Binary ist als Namensschema Code
(Abschnitt~\ref{sec:axis-triad}). Beim Workload reicht der Treiber einen Standard-Workload je Modul durch,
ein Profil bringt seinen eigenen mit; die vorgesehene Voll-Permutation über alle Workloads begegnet der
\emph{Lastprofil-Selektionsverzerrung} --- der Verzerrung eines Vergleichs durch Lastprofile, unter denen
ein
Entwurf bevorzugt gut abschneidet~\cite{blackburn2006dacapo,blackburn2016truth,mytkowicz2009wrongdata} ---;
sie ist nicht implementiert.
\end{nurkurz}

\begin{nurlang}
%% INTERN-272: KA-117
Zu FF2.b (Organ-Achsen): Die widersprüchlichen Literaturbefunde zur optimalen Knoten-/Cache-Line-Größe
--- eine Cache-Line je Knoten bei CSS- und CSB\textsuperscript{+}-Bäumen~\cite{rao1999css,rao2000csb}
gegenüber sechzehn und mehr Cache-Lines bei Hankins und Patel~\cite{hankins2003nodesize} (Knoten von
512 Byte und größer bei 32-Byte-Cache-Lines) sowie mehreren je Knoten bei den
Prefetching-B\textsuperscript{+}-Bäumen von Chen, Gibbons und Mowry~\cite{chen2001prefetch} --- werden nicht aus
ihren Original-Setups übernommen, sondern achsen-isoliert unter sonst gleichen Bedingungen
nachgemessen: Knotentyp, Speicher-Layout und Prefetch sind Organ-Achsen, die einzeln variiert werden,
während System- und Mess-Achsen konstant bleiben; die Zahl ist eine Messung und liegt nicht vor. Zu FF2.c
(Erweiterbarkeit): Ein neuer Prüfling gliedert sich als Lebewesen seiner Gattung ins Genus ein ---
Gattung ist die erste Ebene (Map, Container, Graph), Genus die zweite, und die Genus-Bau-Zulassung
führt die Slot-Sätze je Gattung als übersetzungszeitliches Urteil ---; er dockt am Prüf-Dock an,
durchläuft die Drei-Stufen-Prüfung (Abschnitt~\ref{sec:mess-chain}), und sein Verbund mit der
Cache-Engine führt einen der registrierten Namen: \texttt{Verbund1\_CeOnly} (keine
Prüflings-Direktive), \texttt{Verbund2\_Replace} (die Prüflings-Achse ersetzt die Cache-Engine-Achse
mit Rückfall), \texttt{Verbund2\_Hybrid} (Cache-Engine und Prüfling hybrid je Prüfling; deklariert,
nicht materialisiert --- der Validator sperrt das Token \texttt{merge}, solange der Strategie-Wert nicht implementiert ist)
oder \texttt{Verbund3\_Union} (Vereinigung aller Varianten; vorbereitete Schnittstelle,
Abschnitt~\ref{sec:eval-series}). An der Binary bleibt der Verbund über die hierarchisch erweiterten
Algorithmus-Namen und Flags der Organ-Zeile unterscheidbar --- ein Prüflings-Organ führt sein
Namensraum-Präfix im Eintrag ---; eine eigene Stempel-Zeile braucht es nicht, und ein Versions-Suffix
bezeichnet die Mischung nicht: Das Flag \texttt{e} bezeichnet einen Efficiency-Core-Algorithmus
(Abschnitt~\ref{sec:stamp-model}). Ob der Prüfling eine bestehende Gattung trifft, eine neue
begründet oder eine Mischform darstellt, entscheidet der Gattungs-Schnitt
(Abschnitt~\ref{sec:anatomy-levels}); eine eigene Gattung entstünde erst bei grundlegend anderem
Basis-Achsen-Satz. Der in FF2.c erfragte Betriebsmodus, der die Entwurfsbestandteile genau einer
Veröffentlichung als geschlossenes Achsen-Set isoliert anbietet, ist nicht implementiert und bleibt
Gegenstand des Ausblicks (Abschnitt~\ref{sec:outlook}). Zu FF2.d (Workload-Stand der Technik): Den Stand der Technik der
Workloads für Suchalgorithmen bilden die Schlüssel-Wert-Konvention YCSB mit ihren sechs Lastprofilen
A--F (je ein Operations-Mix über einer Schlüssel-Zugriffsverteilung, von 50/50 Lesen/Aktualisieren
Zipf-verteilt bis zu kurzen Bereichs-Scans und Read-Modify-Write), der Index-Benchmark SOSD, die
TPC-Familie (TPC-C, TPC-DS), die Integration in reale Systeme (SuRF in RocksDB), reale String-Korpora
(URLs, DNA, Wörterbuch-Terme) sowie SPEC~CPU und CloudSuite bei den hardware-nahen Arbeiten
(Abschnitt~\ref{sec:workload-basics-known}); diese Rahmenwerke erzeugen untereinander nicht unmittelbar
vergleichbare Lasten, und jede Veröffentlichung wählt typischerweise das Lastprofil, unter dem ihr
eigener Entwurf am besten abschneidet. Die Arbeit übernimmt daraus YCSB als das registrierte
Last-Rahmenwerk (die Meta-Meta-Achse \texttt{load\_framework} des Mess-Realms) mit der
Laufzeit-Dimension \texttt{workload} über die sechs Token \texttt{YCSB\_A} bis \texttt{YCSB\_F},
geroutet je Permutation aus dem Traversal-Tag des Profils oder seinem
\texttt{<expected\_workload>}-Tag, und fährt sie über reale String-Datensätze mit
Reproduzierbarkeits-Akte (\texttt{url}, \texttt{protein}, \texttt{tpcds-id}, \texttt{trec-terms},
\texttt{dna}, \texttt{xml}; Abschnitt~\ref{sec:eval-workloads}), wobei \texttt{xml} nur eine
Spezifikation ohne Prüfsumme besitzt, weil die Datei nicht lokal vorliegt (n/a); die übrigen
Konventionen sind als Last-Rahmenwerk nicht registriert (n/a). Umsetzungsstand: Der Treiber
fällt für profil-lose Module auf genau einen Workload zurück, die Workload-Option des Laufs; die
Voll-Permutation über alle Lastprofile ist vorgesehen, aber nicht implementiert (Abschnitt~\ref{sec:eval-workloads}). Zur
Plattform-Verifikation: Eine Plattform gilt als verifiziert, wenn das Konformitäts-Gatter je
geladener Tier-Binary die siebzehn in Anhang~\ref{app:interfaces}
markierten \texttt{std::map}-Operationen --- die Operationsmenge seiner zwölf Kern-Prüfblöcke,
darunter \texttt{at}, \texttt{contains}, \texttt{count}, \texttt{insert\_or\_assign}, \texttt{erase},
\texttt{clear}, \texttt{lower\_bound} und \texttt{upper\_bound} --- gegen \texttt{std::map} als Orakel
bestätigt und die Plattform-Probe die gemeldete Achsen-Teilmenge deckt; die fünf Antriebs-Operationen
der Modulgrenze (\texttt{insert}, \texttt{lookup}, \texttt{erase}, \texttt{clear}, \texttt{size}) sind
die ABI-Primitive, aus denen der Host die übrigen Operationen komponiert, nicht der Prüfumfang, und
der Vollausbau auf die volle \texttt{std::map}-Schnittstelle ist nicht implementiert; die
plattform-kalibrierten Schwellwerte liegen nicht vor, ihre Festlegung setzt die Messläufe
voraus. Als \emph{Beleg} dienen Abschnitt~\ref{sec:mess-chain} (Drei-Stufen-Prüfung und
Verbund; Struktur), Abschnitt~\ref{sec:measurement-system} (Mess-Ebenen und Realm; Struktur),
Abschnitt~\ref{sec:anatomy-levels} (Gattungs-Schnitt; Struktur), Anhang~\ref{app:interfaces}
(Vertrags-Teilmenge; Test) sowie für FF2.d Abschnitt~\ref{sec:workload-basics-known}
(Workload-Konventionen; Literatur) und Abschnitt~\ref{sec:eval-workloads} (Workload-Routing und Akten;
Struktur). Die \emph{Folge}: FF2 strukturell beantwortet, FF2.c konzeptionell beantwortet, FF2.d
beantwortet; der Isolations-Modus, die Zahl zu FF2.b und die Voll-Permutation der Workloads sind offen.
\end{nurlang}
\begin{nurkurz}
%% INTERN-272: KA-117
Zu FF2.b (Organ-Achsen): Die widersprüchlichen Literaturbefunde zur optimalen Knoten-/Cache-Line-Größe ---
eine Cache-Line je Knoten~\cite{rao1999css,rao2000csb} gegenüber sechzehn und
mehr~\cite{hankins2003nodesize}
und mehreren je Knoten~\cite{chen2001prefetch} --- werden nicht aus ihren Original-Setups übernommen,
sondern
achsen-isoliert nachgemessen (je eine Organ-Achse variiert, System- und Mess-Achsen konstant); die Zahl
liegt nicht vor. Zu FF2.c (Erweiterbarkeit): Ein neuer Prüfling gliedert sich
als Lebewesen seiner Gattung ins Genus ein (die Genus-Bau-Zulassung urteilt übersetzungszeitlich), dockt am
Prüf-Dock an, durchläuft
die Drei-Stufen-Prüfung (Abschnitt~\ref{sec:mess-chain}) und führt im Verbund mit der Cache-Engine einen der
Namen \texttt{Verbund1\_CeOnly}, \texttt{Verbund2\_Replace}, \texttt{Verbund2\_Hybrid} (deklariert, nicht
materialisiert) oder \texttt{Verbund3\_Union} (vorbereitet, Abschnitt~\ref{sec:eval-series}); an der Binary
bleibt er über die hierarchischen Algorithmus-Namen der Organ-Zeile unterscheidbar, das Flag \texttt{e}
bezeichnet einen Efficiency-Core-Algorithmus, keine Mischung (Abschnitt~\ref{sec:stamp-model}). Ob der
Prüfling eine Gattung trifft, eine neue begründet oder eine Mischform ist, entscheidet der Gattungs-Schnitt
(Abschnitt~\ref{sec:anatomy-levels}; eine eigene Gattung erst bei grundlegend anderem Basis-Achsen-Satz);
der Paper-Isolations-Modus ist nicht implementiert
(Abschnitt~\ref{sec:outlook}). Zu FF2.d (Workload-Stand der Technik): Ihn bilden YCSB mit den sechs
Lastprofilen A--F, der Index-Benchmark SOSD, die TPC-Familie, SuRF in RocksDB, reale String-Korpora sowie
SPEC~CPU und CloudSuite (Abschnitt~\ref{sec:workload-basics-known}). YCSB ist das registrierte
Last-Rahmenwerk (sechs Token, geroutet je Permutation aus dem Profil) über sechs reale String-Datensätze
mit Reproduzierbarkeits-Akte (\texttt{xml} ohne Prüfsumme, n/a;
Abschnitt~\ref{sec:eval-workloads}). Eine Plattform gilt als
verifiziert, wenn das Konformitäts-Gatter je Tier-Binary die siebzehn in Anhang~\ref{app:interfaces}
markierten \texttt{std::map}-Operationen (zwölf Kern-Prüfblöcke, \texttt{std::map} als Orakel) bestätigt
und die Plattform-Probe die gemeldete Achsen-Teilmenge deckt;
die fünf Antriebs-Operationen der Modulgrenze sind ABI-Primitive, nicht der Prüfumfang; Vollausbau und
kalibrierte Schwellwerte fehlen.
\emph{Beleg}: Abschnitte~\ref{sec:mess-chain},~\ref{sec:measurement-system} und~\ref{sec:anatomy-levels}
(Struktur), Anhang~\ref{app:interfaces} (Test); FF2.d: Abschnitt~\ref{sec:workload-basics-known}
(Literatur), Abschnitt~\ref{sec:eval-workloads} (Struktur).
\emph{Folge}: FF2 strukturell, FF2.c konzeptionell, FF2.d beantwortet; Isolations-Modus, die Zahl zu FF2.b
und die Voll-Permutation der Workloads sind offen.
\end{nurkurz}

\begin{nurlang}
\paragraph{FF3 (Prüfling-Vergleich).} Gefragt ist, ob der Prüfling PRT-ART gegenüber den acht
Rang-1-Entwürfen --- ART~\cite{leis2013art}, HOT~\cite{binna2018hot}, Masstree, CoCo-Trie, START,
B\textsuperscript{2}-Baum, Wormhole, SuRF --- und der Rang-2/3-SOTA \emph{nicht} durch asymptotische
Neuheit, sondern durch bessere Mikroarchitektur-Eigenschaften gewinnt --- höhere Cache-Line-Auslastung,
weniger LLC- und dTLB-Misses, geringere Branch-Kosten, kleinere heiße Suchpfade --- bei Punktsuche,
Prefix-Lookup und Prefix-Enumeration, gemessen in den Pflicht-Messreihen A, B und C\@; FF3.a fragt nach
der günstigsten lokalen Seitendarstellung je Verzweigungspunkt, FF3.b nach dem Compile-Time-Zwang.
Der \emph{Status}: FF3 ist offen, FF3.b beantwortet, FF3.a strukturell beantwortet (Vorgabe ohne Messbeleg). Der
Vergleich liegt nicht vor und setzt die ersten Mess-Läufe voraus; der Apparat dafür ist implementiert: das PRT-ART-Profil, die
drei Pflicht-Messreihen --- die Reihe wird mechanisch aus der Verbund-Stufe abgeleitet
(\texttt{Verbund1} und \texttt{Verbund2} liefern Reihe~A, \texttt{Verbund3} Reihe~B, Reihe~C ist
bau-übergreifend; Tabelle~\ref{tab:stage-series}) ---, der \texttt{ExperimentDriver} mit Auto-Pickup
der 33 Profil-XML aus dem Profil-Verzeichnis und der ABI-konforme Adapter mit dem
\texttt{std::map}-Kernvertrag als host-seitiger Kompositions-Schicht über dem Punkt-Such-Primitiv
der Modulgrenze, flankiert von Einfügen und Löschen, mit der geordneten Iteration als additiver
Scan-Fähigkeit (Anhang~\ref{app:interfaces}). Das Konformitäts-Gatter prüft die siebzehn in
Anhang~\ref{app:interfaces} markierten Operationen in zwölf Kern-Prüfblöcken gegen \texttt{std::map}
als Orakel; der Vollausbau auf die volle Schnittstelle ist vorgesehen, aber nicht implementiert. Der
Benachrichtigungs-Haken der Mess-Schicht ist
ein Single-Slot-Hook je Achse, standardmäßig nicht einkompiliert (Opt-in), und kein Teil des Vertrags.
Zu FF3.a: Punktseite --- die Direktindex-Seiten Array$_{256}$ und Array$_{65535}$ --- gegenüber
Kompaktseite --- die Vektor-Seiten Vector$\langle$u8,u8$\rangle$ und Vector$\langle$u16,u16$\rangle$
--- je Verzweigungspunkt und die Umschaltregeln nach Belegungsdichte und Zugriffsart sind Teil des
PRT-ART-Profils; der Prüfling schaltet nach festen Dichte-Schwellen von 25, 50 und 75~Prozent um
(Abschnitt~\ref{sec:prtart-demo}), und die plattform-kalibrierten Regeln sind Vorgaben ohne Messbeleg, solange die
Läufe sie nicht belegen. FF3.b ist bejaht und zur Bauform geworden: Jede Achse gliedert sich in eine
übersetzungszeit-fixierte \emph{Haupt-Achse} und laufzeit-variable \emph{Unter-Achsen}
--- die Haupt-Achsen sind statisch einkompiliert, bei Organ-Achsen als \texttt{binary\_id}-Segment, bei
System-Achsen als Stempel samt Bau-Versionskennung, während Unter-Achsen mit der Stufe
\texttt{runtime} ihre Parameter ohne Neuübersetzung variieren; die Organ-Registry führt die Stufe je
Unter-Achse, die System-Registry die Unter-Achsen-Gruppe \texttt{build\_toolchain} der Komplex-Klammer
(Abschnitt~\ref{sec:measurement-system}). Als \emph{Beleg} dienen Abschnitt~\ref{sec:prtart-demo} (Prüfling
und Seitendarstellung; Struktur), Abschnitt~\ref{sec:sota-instances} (33 Profile; Struktur),
Abschnitt~\ref{sec:measurement-system} (Haupt- und Unter-Achse; Struktur) und
Anhang~\ref{app:interfaces} (Kernvertrag; Test). Die \emph{Folge}: Die Zahl setzt die Reihe~A voraus; die
Hypothesen H1--H3 liefern die Auswerte-Kriterien.
\end{nurlang}
\begin{nurkurz}
\paragraph{FF3 (Prüfling-Vergleich).} Gefragt ist, ob der Prüfling PRT-ART gegenüber den acht
Rang-1-Entwürfen --- ART~\cite{leis2013art}, HOT~\cite{binna2018hot}, Masstree, CoCo-Trie, START,
B\textsuperscript{2}-Baum, Wormhole, SuRF --- und der Rang-2/3-SOTA nicht durch asymptotische Neuheit,
sondern
durch bessere Mikroarchitektur-Eigenschaften in den Pflicht-Messreihen A, B und C gewinnt; FF3.a fragt nach
der günstigsten Seitendarstellung je Verzweigungspunkt, FF3.b nach dem Compile-Time-Zwang. Der
\emph{Status}:
FF3 ist offen, FF3.b beantwortet, FF3.a strukturell beantwortet. Der Vergleich setzt die ersten Mess-Läufe
voraus; der Apparat ist implementiert: das PRT-ART-Profil, die drei Reihen, mechanisch aus der Verbund-Stufe
abgeleitet (Tabelle~\ref{tab:stage-series}), der \texttt{ExperimentDriver} mit Auto-Pickup der 33 Profil-XML
und der ABI-konforme Adapter mit dem \texttt{std::map}-Kernvertrag (Anhang~\ref{app:interfaces}). Das
Konformitäts-Gatter ist
implementiert, sein Vollausbau nicht; der Mess-Hook je Achse ist Opt-in. Zu FF3.a: Punkt- gegenüber
Kompaktseite je Verzweigungspunkt und die Umschaltregeln (feste Dichte-Schwellen 25, 50 und 75~Prozent,
Abschnitt~\ref{sec:prtart-demo}) sind Teil des PRT-ART-Profils, bleiben aber Vorgaben ohne Messbeleg. FF3.b
ist bejaht und Bauform: Jede Achse gliedert sich in eine
übersetzungszeit-fixierte \emph{Haupt-Achse} (\texttt{binary\_id}-Segment oder Stempel) und
laufzeit-variable \emph{Unter-Achsen}, die ihre Parameter ohne Neuübersetzung variieren
(Abschnitt~\ref{sec:measurement-system}).
\emph{Beleg}: Abschnitte~\ref{sec:prtart-demo},~\ref{sec:sota-instances} und~\ref{sec:measurement-system}
(Struktur), Anhang~\ref{app:interfaces} (Test).
\emph{Folge}: Die Zahl setzt die Reihe~A voraus; die Hypothesen H1--H3 liefern die Auswerte-Kriterien.
\end{nurkurz}

\begin{nurlang}
\paragraph{Hypothesen und Zielpfad.} Die Evaluations-Hypothesen H1--H3
(Abschnitt~\ref{sec:eval-hypotheses}) sind vor dem Apparat formuliert worden; der über die Kapitel
erarbeitete Zielpfad --- Dreiteilung der Realms, Stempel-Identität, Mess-Ebenen --- hat sie nicht als
Fundament gelegt, sondern ordnet sie als Prüfsteine der Auswertung ein: H1 (Seitentyp-Kosten) und H3
(Verteilung der Wert-Ablage) werden in Reihe~A je Lastprofil geprüft, H2 (Code-Qualität pro Quelle)
über die tool-berechnete Qualitäts-Akte. Deren Datenbasis ist die gewichtete
\texttt{cppcheck}-Befunddichte je tausend physischer Zeilen über die vendorten Originalquellen; die
Akte führt 33 Einträge, davon 22 als n/a, weil eine Rekonstruktion keine vendorte Originalquelle
besitzt, und PRT-ART besitzt kein Original-Repository (n/a). Die Hypothesenprüfung folgt der Vorgabe,
Kriterium und Schwelle vor der Messung zu nennen~\cite{wohlin2012experimentation}, und die
Wiederholungszahl je Zelle richtet sich nach der geforderten Konfidenz~\cite{kalibera2013rigorous}.
%% INTERN-272: KA-118
H2 gilt als angenommen, wenn über die Vollprofile mit vendorter Originalquelle die Rangkorrelation nach
Spearman zwischen Befunddichte und gemessenem Durchsatz je Lastprofil und Zelle (Plattform, Workload,
Regime) negativ ist mit $|\rho| \geq 0{,}5$ bei $\alpha = 0{,}05$; als abgelehnt bei $|\rho| < 0{,}3$;
dazwischen als nicht entschieden; unterhalb von acht Profilen mit Score ist H2 nicht prüfbar. Das
Hand-Audit über sieben Bewertungs-Achsen bleibt die qualitative Einordnung. Die
PRT-ART-Telemetrie-Metriken H1--H3 sind Mess-Organe des Prüflings, nicht diese Hypothesen
(Abschnitt~\ref{sec:eval-hypotheses}); und die Variante ChainRef der Wert-Ablage ist Gegenstand von H3, sie ist kein
Abbruchkriterium für H2.
\end{nurlang}
\begin{nurkurz}
\paragraph{Hypothesen und Zielpfad.} Die Evaluations-Hypothesen H1--H3 (Abschnitt~\ref{sec:eval-hypotheses})
sind vor dem Apparat formuliert; der Zielpfad ordnet sie als Prüfsteine der Auswertung ein: H1
(Seitentyp-Kosten) und H3 (Verteilung der
Wert-Ablage) werden in Reihe~A je Lastprofil geprüft, H2 (Code-Qualität pro Quelle) über die tool-berechnete
\texttt{cppcheck}-Akte mit 33 Einträgen, davon 22 als n/a ohne vendorte Originalquelle; PRT-ART besitzt kein
Original-Repository (n/a). Kriterium und Schwelle stehen vor der Messung~\cite{wohlin2012experimentation},
die Wiederholungszahl je Zelle richtet sich nach der geforderten Konfidenz~\cite{kalibera2013rigorous}.
%% INTERN-272: KA-118
H2 gilt als angenommen, wenn über die Vollprofile mit vendorter Originalquelle die Rangkorrelation nach
Spearman zwischen Befunddichte und gemessenem Durchsatz je Lastprofil und Zelle (Plattform, Workload,
Regime) negativ ist mit $|\rho| \geq 0{,}5$ bei $\alpha = 0{,}05$; als abgelehnt bei $|\rho| < 0{,}3$;
dazwischen als nicht entschieden; unterhalb von acht Profilen mit Score ist H2 nicht prüfbar. Das
Hand-Audit über sieben Bewertungs-Achsen bleibt die qualitative Einordnung. Die PRT-ART-Telemetrie-Metriken
H1--H3 sind Mess-Organe des Prüflings, nicht diese Hypothesen (Abschnitt~\ref{sec:eval-hypotheses}); die
Variante ChainRef der Wert-Ablage ist Gegenstand von H3, kein Abbruchkriterium für H2.
\end{nurkurz}

\begin{nurlang}
\paragraph{FF4 (Compile-Time vs.\ Laufzeit).} Gefragt ist, welche Achsen-Eigenschaften zur
Übersetzungszeit fixiert werden müssen, um plattform-optimierte, provenienz-nachweisbare Binaries je
Permutation zu erzeugen, und welche zur Laufzeit gewählt werden können bzw.\ müssen (bit-codierte
Binary-Identität je Permutation, laufzeit-adaptive Layout-Wahl), ohne die Vergleichbarkeit (gleiche
Compiler- und Flag-Basis) zu verletzen --- und ob es durch die Messung möglich wird, ISA+OS-optimierte
Binaries für eine Suchalgorithmus-Rekombination zu erzeugen, deren heuristische Optimierung perfekt
auf ihre Laufzeitumgebung zugeschnitten ist; FF4.a fragt, ob sich der gesamte Ablauf --- vom
plattformbezogenen Bau aller Binary-Permutationen über den Messlauf und die Auswahl der besten
gemessenen Binary oder Binary-Rekombination bis zur Auswertung in CSV-/xlsx-Dateien, LaTeX und
Diagrammen --- durchgehend aus einer XML-Beschreibung des Experiments erzeugen lässt; die Erweiterung
FF4.E fragt, ob die Grenze entlang der Achsen-Kategorien verläuft.
Der \emph{Status}: FF4 ist strukturell beantwortet, FF4.a strukturell beantwortet, FF4.E beantwortet. Die Grenze verläuft entlang der
Kategorien: Die drei System-Haupt-Achsen \texttt{target\_isa}, \texttt{operating\_system} und
\texttt{external\_utils} samt ihrer Komplex-Klammer \texttt{build\_target\_complex}
(Abschnitt~\ref{sec:complex-axis}; keine vierte Haupt-Achse) gehören zur Übersetzungszeit-Identität
der Binary; sie erscheinen als System-Stempel-Zeile und Bau-Versionskennung, nie als
\texttt{binary\_id}-Segment. Die Provenienz sichert der Tier-Fingerprint: ein SHA-512 über elf Glieder
(Fingerprint-Format~6), darunter das Overlay-Glied als ein SHA-512 über die konkatenierten Quelldateien
des Overlay-Schnitts --- je Achse die Dateimenge ihrer eigenen Implementierung plus die gemeinsame
Anatomie-Substanz ---, sodass eine reine Quelltext-Änderung eine neue Identität ergibt, und das
Toolchain-Glied mit der Bau-Compiler-Kennung und der Optimierungs-Identität; der Planer-Fingerprint ist
ein SHA-256 über sein eigenes Preimage. Die Abfrage \texttt{get\_compiler()} einer Achse liefert
dagegen den Original-Compiler der Veröffentlichung aus dem Manifest der vendorten Quelle --- eine
Herkunftsangabe, keine Bau-Kennung. Mess-Wahlen bleiben laufzeit-variabel; Layout- und Seiten-Selektion
bleiben über die mixed-radix-Ordnung der \texttt{binary\_id} (die statische Binary-Sicht) und die
laufzeit-dynamische Cache-Line-Politik wählbar, ohne die gemeinsame Compiler- und Flag-Basis zu
verletzen. Wo eine Achse beide Seiten der Grenze braucht, ist sie zweigeteilt: Last-Rahmenwerk,
Telemetrie und Hardware-Erkennung laufen über die Laufzeit des Planers \emph{und} die Übersetzungszeit
des Builders --- die Unter-Achsen-Freigabe der Vorstufe wird zur Haupt-Achsen-Annahme der Folgestufe
(Abschnitt~\ref{sec:axis-triad}) ---, und die Hardware-Erkennung setzt die Übersetzungszeit-Stempel der
Tier-Binaries aus den zur Builder-Laufzeit erhobenen Ist-Eigenschaften (Abschnitt~\ref{sec:hw-probe}).
Gestempelt wird stets die Plattform-\emph{Klasse}, nie ein Messwert; der \emph{Passungs-Stempel}, der
die gestempelte Klasse bei jedem Mess-Lauf gegen das Laufzeit-Verdikt der Hardware-Erkennung prüft, ist
entworfen, aber nicht implementiert (Abschnitt~\ref{sec:outlook}). Die vorwärtsgerichtete Teilfrage, ob die Messung
selbst ISA- und betriebssystem-optimierte Binaries für eine Rekombination zuschneiden kann, greift der
Ausblick als Heuristik-Automation auf; die Betriebssystem-Dimension ist mit drei deklarierten Werten
der System-Achse (\texttt{linux}, \texttt{windows}, \texttt{macos}, je mit eigener Probe) und den zwei
Mess-Regimes im Messplan angelegt --- das zweite Regime, ein immutables Betriebssystem, ist entworfen,
nicht umgesetzt (Abschnitt~\ref{sec:eval-platforms}). Zu FF4.a (deklarative Kette): Die Kette ist
durchgehend deklarativ angelegt und als zwei Binaries gebaut. Der Planer
\texttt{comdare-experiment-planner} prüft die Anwender-XML rein lesend gegen die einkompilierten
Registry-Kataloge und emittiert aus dem einen Director-Walk den deterministischen Plan --- zwei Läufe
sind byte-gleich, der Plan-Kopf annotiert das Registry-Trio, und jeder Bau-Schritt weist den tatsächlich
erzeugten Datei-Stamm aus derselben Orchestrator-Funktion aus, die auch der Bau aufruft ---; der
Mess-Treiber \texttt{comdare-messung-driver} emittiert die System-Permutationen und
Tier-Bau-Aufträge seines freigegebenen Bau-Raums und fährt die Sieben-Phasen-Mess-Pipeline von
\textsc{Enumerate} bis \textsc{Persist} (CSV und JSON); die nachgelagerte Neun-Stufen-Werkzeugkette
rendert daraus die xlsx-Mappe, die booktabs-Tabellen, die TikZ-Diagramme und die Anhang-Dateien, die
das Manuskript einbindet --- eine XML-Quelle, ein Bau-Kanal, kein Schritt von Hand
(Abschnitt~\ref{sec:impl-pipeline}, Abbildung~\ref{fig:eval-xml-kette}). Die Auswahl der besten
Binary oder Rekombination ist als Arbeitsmodus \texttt{compare} der Ablauf-Kette registriert und
validierbar; sein modus-spezifischer Ablauf und der \texttt{release} sind vorgesehen, aber nicht
implementiert --- beide sind wählbare, validierbare Etiketten mit \texttt{release}-gleicher Bausemantik
(Abschnitt~\ref{sec:anatomy-levels}).
Umsetzungsstand: Die Kette ist in einem ersten durchgehenden Lauf nachgewiesen --- die Smoke-Messreihe über 43
Permutationen im Anhang ist ihr Erzeugnis ---; der headless Dienst, der XML-Erstellungsmodus, die
\texttt{<publish>}-Sektion und die CSV-Baumform sind vorgesehen, aber nicht implementiert. Als \emph{Beleg} dienen
Abschnitt~\ref{sec:complex-axis} (Komplex-Klammer; Struktur), Abschnitt~\ref{sec:hw-probe}
(zweigeteilte Erkennung; Struktur), Abschnitt~\ref{sec:stamp-model} (Stempel-Zeilen; Struktur),
Abschnitt~\ref{sec:impl-system-axes} (Rückkehr-Sperren und Ordnungs-Prüfungen; Test) sowie für FF4.a
Abschnitt~\ref{sec:impl-pipeline} (Planer, Treiber und Mess-Pipeline; Struktur) und
Abschnitt~\ref{sec:mess-chain} (Kette von der XML zur Tier-Binary; Struktur). Die \emph{Folge}: FF4
und FF4.a strukturell beantwortet; Heuristik-Automation, Passungs-Stempel sowie der Ablauf von
\texttt{compare} und \texttt{release} offen.
\end{nurlang}
\begin{nurkurz}
\paragraph{FF4 (Compile-Time vs.\ Laufzeit).} Gefragt ist, welche Achsen-Eigenschaften zur Übersetzungszeit
fixiert werden müssen, um plattform-optimierte, provenienz-nachweisbare Binaries je Permutation zu erzeugen,
und welche zur Laufzeit wählbar bleiben, ohne die gleiche Compiler- und Flag-Basis zu verletzen; FF4.a fragt
nach der durchgehenden XML-Kette vom Bau bis zur Auswertung, FF4.E, ob die Grenze entlang der
Achsen-Kategorien verläuft. Der \emph{Status}: FF4 und FF4.a sind strukturell beantwortet, FF4.E
beantwortet.
Die Grenze verläuft entlang der Kategorien: Die drei System-Haupt-Achsen samt Komplex-Klammer
\texttt{build\_target\_complex} (Abschnitt~\ref{sec:complex-axis}) gehören zur Übersetzungszeit-Identität
der Binary und erscheinen als System-Stempel-Zeile, nie als \texttt{binary\_id}-Segment. Die Provenienz
sichert der Tier-Fingerprint (SHA-512, elf Glieder, Format~6; das Overlay-Glied hasht die Quelldateien
des Overlay-Schnitts --- je Achse die Dateien ihrer eigenen Implementierung plus die gemeinsame
Anatomie-Substanz ---, das Toolchain-Glied trägt Bau-Compiler-Kennung und Optimierungs-Identität); der
Planer-Fingerprint ist ein SHA-256,
und \texttt{get\_compiler()}
liefert den Original-Compiler der Veröffentlichung --- Herkunftsangabe, keine Bau-Kennung. Mess-Wahlen
bleiben laufzeit-variabel, Layout- und Seiten-Selektion über statische \texttt{binary\_id} und dynamische
Cache-Line-Politik wählbar.
Zweigeteilt (Planer-Laufzeit und Builder-Übersetzungszeit) sind Last-Rahmenwerk, Telemetrie und
Hardware-Erkennung (Abschnitte~\ref{sec:axis-triad} und~\ref{sec:hw-probe}).
Gestempelt wird stets die Plattform-\emph{Klasse}, nie ein Messwert; der \emph{Passungs-Stempel}, der die
gestempelte Klasse bei jedem Mess-Lauf gegen das Laufzeit-Verdikt der Hardware-Erkennung prüft, ist
entworfen, nicht implementiert (Abschnitt~\ref{sec:outlook}); die vorwärtsgerichtete Teilfrage, ob die
Messung selbst ISA- und betriebssystem-optimierte Binaries für eine Rekombination zuschneiden kann, greift
der Ausblick als Heuristik-Automation auf; die Betriebssystem-Dimension ist mit drei
Werten und zwei Mess-Regimes angelegt, das zweite Regime nicht umgesetzt
(Abschnitt~\ref{sec:eval-platforms}).
Zu FF4.a: Die Kette ist durchgehend deklarativ und als zwei Binaries gebaut: Der Planer prüft die
Anwender-XML rein lesend und emittiert deterministisch den Plan, der Mess-Treiber fährt die
Sieben-Phasen-Mess-Pipeline von \textsc{Enumerate} bis \textsc{Persist}, die Neun-Stufen-Werkzeugkette
führt bis zum Manuskript (Abschnitt~\ref{sec:impl-pipeline}, Abbildung~\ref{fig:eval-xml-kette}).
\texttt{compare} (Auswahl der besten Binary) und \texttt{release} sind als Arbeitsmodi registriert, ihr
Ablauf ist nicht implementiert (Abschnitt~\ref{sec:anatomy-levels}); Umsetzungsstand: ein erster
durchgehender Lauf (43 Permutationen) liegt vor; headless Dienst, XML-Erstellungsmodus,
\texttt{<publish>}-Sektion und CSV-Baumform fehlen.
\emph{Beleg}: Abschnitte~\ref{sec:complex-axis},~\ref{sec:hw-probe} und~\ref{sec:stamp-model} (Struktur),
Abschnitt~\ref{sec:impl-system-axes} (Test); FF4.a: Abschnitte~\ref{sec:impl-pipeline}
und~\ref{sec:mess-chain} (Struktur).
\emph{Folge}: FF4 und FF4.a strukturell beantwortet; Heuristik-Automation, Passungs-Stempel sowie der Ablauf
von \texttt{compare} und \texttt{release} offen.
\end{nurkurz}

\begin{table}[htbp]\centering\footnotesize
\begin{tabular}{@{}>{\raggedright\arraybackslash}p{0.9cm}>{\raggedright\arraybackslash}p{3.0cm}>{\raggedright\arraybackslash}p{2.0cm}>{\raggedright\arraybackslash}p{4.8cm}>{\raggedright\arraybackslash}p{3.3cm}@{}}
\toprule
Frage & Teilfrage & Status & Beweis (Abschnitte; Belegart) & Rest \\
\midrule
FF0 & Hauptfrage: aktiv, messgetrieben gegen passiv, statisch & offen & {\ref{sec:axis-triad}} (Struktur), {\ref{sec:impl-system-axes}} (Test), {\ref{sec:eval-method}} (Messplan) & Zahl aus Reihe~A nach dem vollständigen Bau \\
FF0.E & Zuordenbarkeit eines Gewinns zu Organ, System, Mess & strukturell beantwortet & {\ref{sec:axis-triad}}, {\ref{sec:stamp-model}} (Struktur) & Messung, die die Zuordnung ausübt \\
FF1 & kanonische Zerlegung, Permutationsmatrix, Profile & strukturell beantwortet & {\ref{sec:axis-triad}}, {\ref{sec:mess-chain}} (Struktur), {\ref{sec:impl-contracts}} (Test) & Lauf über alle achtzehn Achsen; Profil-Anschluss der Meta-Meta-Achse \\
FF1.a & Kategorien-Schnitt der Zerlegung & beantwortet & {\ref{sec:axis-triad}} (Struktur), {\ref{sec:impl-system-axes}} (Test) & --- \\
FF1.b & Organ-Kopplungen, ISA-/SIMD-Variation & offen (Messung) & {\ref{sec:axis-triad}} (Struktur) & Cache-Line-Füllverhalten je System-Permutation \\
FF2 & Mess-Pipeline, drei Reihen, drei Ebenen, Verifikation & strukturell beantwortet & {\ref{sec:mess-chain}}, {\ref{sec:measurement-system}} (Struktur) & Voll-Permutation der Workloads; kalibrierte Schwellwerte \\
FF2.a & Mess-Achsen als eigener Realm & beantwortet & {\ref{sec:axis-triad}} (Struktur), {\ref{sec:impl-system-axes}} (Test) & Katalog der Erfolgs-Parameter je Achse \\
FF2.b & Knoten-/Cache-Line-Größe achsen-isoliert & strukturell beantwortet & {\ref{sec:measurement-system}} (Struktur) & Messung über Knotentyp, Layout, Prefetch \\
FF2.c & Erweiterbarkeit, Eingliederung neuer Prüflinge & beantwortet (konzeptionell) & {\ref{sec:anatomy-levels}}, {\ref{sec:mess-chain}} (Struktur), {\ref{app:interfaces}} (Test) & Paper-Isolations-Modus; Materialisierung von \texttt{Verbund2\_Hybrid} \\
FF2.d & Workload-Stand der Technik & beantwortet & {\ref{sec:workload-basics-known}} (Literatur), {\ref{sec:eval-workloads}} (Struktur) & Voll-Permutation der Workloads; nur YCSB als Last-Rahmenwerk registriert \\
FF3 & PRT-ART gegen SOTA in den Reihen A--C & offen & {\ref{sec:prtart-demo}}, {\ref{sec:sota-instances}} (Struktur) & Reihe~A\@; Vollausbau des Vertrags \\
FF3.a & Seitendarstellung, Umschaltregeln & strukturell beantwortet (Vorgabe ohne Messbeleg) & {\ref{sec:prtart-demo}} (Struktur) & plattform-kalibrierte Regeln aus den Läufen \\
FF3.b & Compile-Time-Zwang & beantwortet & {\ref{sec:measurement-system}} (Struktur), {\ref{sec:impl-contracts}} (Test) & --- \\
FF4 & Übersetzungszeit gegen Laufzeit, Provenienz & strukturell beantwortet & {\ref{sec:stamp-model}}, {\ref{sec:hw-probe}} (Struktur), {\ref{sec:impl-system-axes}} (Test) & Passungs-Stempel; Heuristik-Automation \\
FF4.a & deklarative Kette von der XML bis zur Auswertung & strukturell beantwortet & {\ref{sec:impl-pipeline}}, {\ref{sec:mess-chain}} (Struktur) & Ablauf von \texttt{compare} und \texttt{release}; Dienst, Erstellungsmodus, CSV-Baum \\
FF4.E & Grenze entlang der Achsen-Kategorien & beantwortet & {\ref{sec:complex-axis}}, {\ref{sec:stamp-model}} (Struktur) & --- \\
\bottomrule
\end{tabular}
\caption{Die Beweiskette je Haupt- und Teilfrage: Status in drei Stufen (beantwortet, strukturell
beantwortet, offen), die stützenden Abschnitte mit ihrer Belegart (Struktur, Test, Messung; für FF2.d der
Literaturbeleg) und der offene Rest ohne Zahl.}\label{tab:ff-beweiskette}
\end{table}

\begin{nurlang}
%% INTERN-272: KA-119
\paragraph{Umsetzungsstand.} Der deklarierte Bau-Umfang --- je System-Permutation die 131072
Organ-Permutationen des golden-Referenzkatalogs; der Katalog-Kommentar schätzt rund 224~GB und rund
24~Stunden je vollständigem Bau, eine Schätzung, kein Messwert --- ist nicht durchgeführt; die End-to-End-Läufe des
Mess-Treibers über die Zielplattformen (x86-64 mit AVX2 und AVX-512, aarch64 mit NEON und
SVE\footnote{NEON und SVE\@: die SIMD-Erweiterungen der ARM-Architektur --- NEON mit fester
128-Bit-Breite, SVE (Scalable Vector Extension) mit implementierungsabhängiger Vektorbreite.}) liegen
nicht vor, ebenso die Abgabe-Messung über den 320er-Katalog. Belegt mit Lauf sind zwei Teil-Durchführungen:
der erste durchgehende Lauf der Kette von der Experiment-XML über Bau und Messung bis zum eingebundenen Anhang
(Abschnitt~\ref{sec:eval-xml}) --- ein Mess-Job der Smoke-Stufe auf der AMD-Zen-5-Maschine\footnote{Provenienz
des Smoke-Laufs: Pipeline 15800, Mess-Job 377503, Host prod1 (Anhang~\ref{app:measurements}); ein
vorausgehender erster durchgehender Lauf derselben Kette (Pipeline 15764, Job 376333) ergab bei 10000 Operationen
1199,047~ns je Operation, seine messende Maschine ist nicht dokumentiert und die Differenz nicht
bewertet.} mit 10000 Operationen und 671,5~ns je Operation (p50 der
\texttt{lookup}-Operation 750~ns) als realer Messwert des Weges --- und der Kalibrierlauf auf der
AMD-Maschine, der 250 von 250
Tier-Binaries gebaut und den Bau-Takt auf 0,521~s je Binary sowie den Platzbedarf auf rund 457~kB je
Binary kalibriert hat; daraus folgt ein Lagerbedarf von rund 240~GB für den vollständigen Bau von 524288
Tier-Binaries über die vier System-Permutationen des golden-Profils (zwei Optimierungsstufen mal zwei
SIMD-Optionen). Die Smoke-Reihe --- 43 Permutationen in zwei Teilblöcken: 27 kartesische
SIMD$\times$Layout$\times$Allokator-Kombinationen plus 16 über Knotenbreite, Pfadkompression, interne
Suche und Telemetrie --- liegt in Anhang~\ref{app:measurements} vor; ihre Provenienz ist der Datenstand,
unter dem sie erhoben wurde: sieben variierte Achsen jener Zählung, von denen SIMD und Telemetrie in der
kanonischen Ordnung keine Organ-Achsen mehr sind. Sie approximiert Achsen-Effekte teils über
Wall-Clock-Proxys, und einzelne Effekte sind Instrumentierungs-Artefakte\footnote{Instrumentierungs-Artefakt
(probe effect): die Veränderung des beobachteten Verhaltens durch das Mess-Instrument
selbst~\cite{gait1986probe}.} statt Organ-Effekte; die Vorbehalte je Reihe führt
Tabelle~\ref{tab:le:limitierung} (Abschnitt~\ref{sec:measurements:limitations}). Die apparat-seitige
Abnahme je Paket läuft über das Zwei-Gate-Protokoll (Abschnitt~\ref{sec:impl-contracts}). Die \emph{Folge}:
FF0 und FF3 bleiben offen, solange Reihe~A nicht gelaufen ist; alle achtzehn Achsen sind für die Variation vorgesehen (nicht gemessen).
\end{nurlang}
\begin{nurkurz}
%% INTERN-272: KA-119
\paragraph{Umsetzungsstand.} Der deklarierte Bau-Umfang (je System-Permutation die 131072
Organ-Permutationen des golden-Referenzkatalogs; der Katalog-Kommentar schätzt rund 224~GB und 24~Stunden je
Bau, kein Messwert) ist nicht durchgeführt; die End-to-End-Läufe über die Zielplattformen (x86-64, aarch64)
liegen nicht vor, ebenso die Abgabe-Messung über den 320er-Katalog.
Belegt sind zwei Teil-Durchführungen: der erste durchgehende Lauf der Kette (Abschnitt~\ref{sec:eval-xml})
--- ein Mess-Job der Smoke-Stufe auf der AMD-Zen-5-Maschine (Pipeline 15800, Mess-Job 377503, Host prod1)
mit 10000 Operationen und 671,5~ns je Operation (p50 \texttt{lookup} 750~ns; ein Vorlauf, Pipeline 15764,
ergab 1199,047~ns auf nicht dokumentierter Maschine) --- und der Kalibrierlauf (250 von 250 Tier-Binaries,
0,521~s und rund 457~kB je Binary; daraus rund 240~GB für 524288 Tier-Binaries über die vier
System-Permutationen des golden-Profils). Die Smoke-Reihe (43 = 27 + 16 Permutationen) liegt in
Anhang~\ref{app:measurements} vor; ihr Datenstand zählte SIMD und Telemetrie noch als Organ-Achsen. Sie
approximiert Achsen-Effekte teils über Wall-Clock-Proxys, einzelne Effekte sind
Instrumentierungs-Artefakte (Veränderung des Verhaltens durch das Mess-Instrument
selbst~\cite{gait1986probe}); die Vorbehalte je Reihe führt
Tabelle~\ref{tab:le:limitierung} (Abschnitt~\ref{sec:measurements:limitations}), die Abnahme je Paket läuft
über das Zwei-Gate-Protokoll (Abschnitt~\ref{sec:impl-contracts}). Die \emph{Folge}: FF0 und
FF3 bleiben offen, solange Reihe~A nicht gelaufen ist; alle achtzehn Achsen sind für die Variation
vorgesehen (nicht gemessen).
\end{nurkurz}

\begin{nurlang}
%% INTERN-272: KA-120
\section{Limitierungen}\label{sec:limitations}
Die folgenden Grenzen beschreiben, was ein Anwender mit dem System nicht erreichen kann --- abgesehen
von den natürlichen und selbst gewählten Grenzen der Maschine, auf der er arbeitet; Vorbehalte
einzelner Messreihen führt der Mess-Anhang (Tabelle~\ref{tab:le:limitierung}), Implementierungsstände
das Kapitel Implementierung (Kapitel~\ref{ch:impl}), und der Umsetzungsstand von vollständigem Bau und
Smoke-Reihe steht in Abschnitt~\ref{sec:answers}. Zwei Grenzen folgen aus dem Lager-Prinzip, die
übrigen sind aus der Architektur abgeleitet: Maschinenspezifische empirische Optimierung ist der
Wesenszug cache-bewusster Verfahren --- ATLAS~\cite{whaley2001atlas} und FFTW~\cite{frigo1998fftw}
vermessen die Maschine, bevor sie den Code wählen ---, und cache-oblivious Verfahren~\cite{frigo1999cache}
sind der Gegenpol ohne Maschinenbindung; ein Apparat, der die Maschinenbindung zum Prinzip macht, erbt
ihre Kosten. Die Modul-Bodies sind keine Grenze: Der Codegenerator des Mess-Pfads emittiert je
Permutation die achtzehn voll qualifizierten Achsen-Typen samt Umbrella-Include, der Alt-Pfad mit dem
Prüflings-Stub ist durch eine strikte Sperre vom Mess-Pfad getrennt, und die Nullaussagen fest zugeordneter
Pool-Familien sind explizit kodierte Nicht-Erhebung mit Grund (Abschnitt~\ref{sec:impl-pipeline};
Tabelle~\ref{tab:le:limitierung}).
\begin{itemize}
  \item \textbf{Versionswechsel eines Achsen-Algorithmus:} Ändert der Anwender die Version eines
  Achsen-Algorithmus, wird jede Tier-Binary, deren Overlay-Schnitt die Datei enthält, zu einer neuen
  Identität --- das Overlay-Glied des Fingerprints hasht alle Quellen des Schnitts, und der Versions-Lock
  verlangt für jede Änderung an einer verzeichneten Datei den bewussten Regenerations-Commit
  (Abschnitt~\ref{sec:impl-contracts}); der Lager-Bestand dieser Binaries samt Messung und Auswertung
  ist neu zu erzeugen, denn der Fingerprint des CacheEngineBuilders ist Provenienz, kein
  Wiederverwendungs-Kriterium. Was der Anwender nicht kann: einen einzelnen Algorithmus austauschen und
  die Messungen der übrigen Binaries behalten.
  \item \textbf{Lager je Maschinentyp und Hardware-Konfiguration:} Messwerte sind an die
  Hardware-Identität der messenden Maschine geschlüsselt --- Bestand~2 des Lagers ist über den
  Fingerprint der vermessenen Binary und die Hardware-Identität adressiert, die Plattform-Zelle daneben
  geklammert (Abschnitt~\ref{sec:lager}) ---; das gesamte Lager ist je Maschinentyp und
  Hardware-Konfiguration einzeln zu messen. Was der Anwender nicht kann: Messwerte einer Maschine auf
  eine andere übertragen. Das ist der Preis der maschinenspezifischen Optimierung
  (ATLAS~\cite{whaley2001atlas}, FFTW~\cite{frigo1998fftw}) und der Grund für die drei
  Maschinen-Signaturen der System-Registry.
%% INTERN-272: KA-121
  \item \textbf{Deckung der Hardware-Zähler je Mikroarchitektur:} Im Mess-Pfad ist die
  Zähler-Erhebung aktiv --- die emittierten Mess-Jobs setzen den PMC-Schalter, nur der Test-Bau baut
  ohne Hardware-Zugriff. Erhoben werden vier generische Zähler (L1-, Last-Level-, dTLB- und
  Branch-Misses über \texttt{perf\_event\_open}\footnote{\texttt{perf\_event\_open}(2): der
  Linux-Systemaufruf, der einen Hardware- oder Software-Zähler als Dateideskriptor öffnet; generische
  Ereignisse bildet der Kernel je Mikroarchitektur ab, modellspezifische verlangen eine rohe Kodierung.})
  und zwei modell-gebundene Rohzähler (L2-Demand-Misses und Fills aus dem Cache eines fremden CCX als
  Kohärenz-Beobachtbare), Letztere nur, wo die Kodierung durch Messung gegengeprüft ist --- AMD Zen~5; ohne
  Katalog-Eintrag führen die Felder eine explizit kodierte
  Nicht-Erhebung\footnote{Explizit kodierte Nicht-Erhebung: Der fehlende Wert ist ein eigener Zustand
  (n/a mit Grund), nie eine Null; die Statistik spricht von geplant fehlenden Daten, deren
  Fehlmechanismus bekannt ist und darum keine Verzerrung einträgt~\cite{graham2006planned,graham2009missing}.}
  (n/a), nie eine Null, und für Intel gibt es keine Rohbelegung, bis ein Intel-Host die Gegenprüfung
  liefert. Das Last-Level-Ereignis scheitert auf Zen~5 beim Öffnen (keine generische Abbildung),
  RAPL\footnote{RAPL (Running Average Power Limit): die Energie-Zähler der Prozessor-Pakete, unter
  Linux über die powercap-Sysfs-Zone lesbar; die Zone ist oft nur mit Root-Rechten lesbar.} ist
  best-effort und ohne Zonen-Leserecht leer, und IPC ist als Kategorie deklariert, nicht erhoben. In
  der Smoke-Reihe sind die Hardware-Spalten nicht erhoben (n/a); die tabellierten Nullen der
  Anhang-Tabellen sind Darstellung, nicht Messung (Tabelle~\ref{tab:le:limitierung}, Zeile~1). Was der
  Anwender nicht kann: L2- und Kohärenz-Zähler auf anderen Modellen als Zen~5 erhalten
  (Abbildung~\ref{fig:measurement-triangle}).
  \item \textbf{Hardware-Passung ohne automatische Prüfung:} Die Hardware-Erkennung liegt in drei
  Funktionen vor: der Erhebung (die ISA-Komplex-Achse mit SPD-/XMP-Deklaration der Speichertaktung und
  die Planer-Probe, die je Lauf den Host nach seinen Zählern fragt), der Übersetzungszeit-Factory (die
  Tier-Stempel aus den Ist-Eigenschaften) und dem Laufzeit-Verdikt; die Maschinen-Deklarationen der
  System-Registry sind Deklarationen mit konservierter Quelle (\texttt{dmidecode}). Nicht vorhanden
  sind die Anbindung an Runner und Prüf-Dock, der additive Ergebnis-Rückschrieb, der
  Provenienz-Roundtrip der Deklarationen und der Passungs-Stempel. Was der Anwender nicht kann: die
  Passung zwischen Maschine und Binary automatisch prüfen lassen (Abschnitte~\ref{sec:hw-probe}
  und~\ref{sec:impl-system-axes}).
  \item \textbf{Ausgabe ohne Blatt-Auffächerung:} Die Ergebnis-Mappe entsteht bedingungslos im
  Speicher; xlsx ist der Standard, die CSV-Form ihr wählbares, immer aus ihr erzeugtes Gegenstück --- kein
  einzelnes Kind-Dokument, sondern ein geordneter Baum aus CSV-Dateien mit einem Ordner je Blatt ---, beide
  zusammen sind wählbar
  (Abschnitt~\ref{sec:eval-identity}). Die Mappe besitzt einen konstanten Blattschlüssel --- alle Zeilen
  in einem Blatt ---; die Auffächerung in ein Blatt je Unter-Achsen-Permutation plus Info-Blatt ist als
  Filterkette der Auswertung vorgesehen, aber nicht implementiert; die CSV-Projektion ist eine flache Datei,
  die Baumform nicht implementiert. Was der Anwender nicht kann: Ergebnisse je
  Unter-Achsen-Permutation als eigene Blätter erhalten oder die CSV-Ablage als Baum erhalten.
  \item \textbf{Permutationsraum und Budget:} Ein vollständiger Bau des golden-Referenzkatalogs kostet je
  System-Permutation nach der Schätzung des Katalogs rund 224~GB und rund 24~Stunden, nach der
  Kalibrierung rund 457~kB und 0,521~s je Tier-Binary (Abschnitt~\ref{sec:answers}); der Planer rechnet
  die Mächtigkeit je Profil und bietet im Subkommando \texttt{simulate} Budget-Eingänge (Bytes je
  Binary, Lager-Budget, Mess-Teilmenge, Kandidaten). Was der Anwender nicht kann: den vollen Produktraum
  aller Achsen-Kataloge materialisieren --- er wählt Kataloge oder Teilmengen, und die Wiederholungszahl
  je Messung bleibt eine Abwägung von Aufwand und Konfidenz~\cite{kalibera2013rigorous}.
  \item \textbf{Kein Laufzeit-Wechsel zwischen Tier-Binaries:} Die Hybrid-Stufe ist als
  Klassifikation, Dock-Schicht, XML-Parser und Binary-Proxy implementiert
  (Abschnitt~\ref{sec:impl-hybrid-lager}); der \texttt{.so}-Export der Hybrid-Stufe, die Eviction und
  der Router sind nicht implementiert --- Reroute und Router setzen Messkurven voraus. Was der Anwender nicht
  kann: ein Hybrid-Tier betreiben, das zur Laufzeit zwischen vorgeladenen Binaries wechselt; auch die
  Zustands-Übergabe beim Wechsel --- der Memento-Abgleich mit den zwei Release-Artefakten single und
  hybrid im Lager --- ist entworfen, nicht implementiert.
  \item \textbf{Betriebssystem-Regime:} Die System-Achse \texttt{operating\_system} ist mit
  \texttt{linux}, \texttt{windows} und \texttt{macos} deklariert und besitzt je eine Probe; gemessen wird
  unter root-Linux mit vollem Zähler-Zugriff, das zweite Regime (immutables Betriebssystem) ist
  entworfen, nicht umgesetzt (Abschnitt~\ref{sec:eval-platforms}), und ein Windows-Zählertreiber liegt
  nicht vor (Tabelle~\ref{tab:le:limitierung}, Zeile~1). Was der Anwender nicht kann: Messreihen unter
  Windows oder macOS oder ohne perf-Rechte (\texttt{perf\_event\_paranoid}, RAPL-Zonen) erheben.
  \item \textbf{Keine isolierte Nachmessung einer einzelnen Veröffentlichung:} Ein Betriebsmodus, der
  die Entwurfsbestandteile genau einer Veröffentlichung als geschlossenes Achsen-Set anbietet, ist nicht
  implementiert; je Veröffentlichung existieren das Profil-XML (33) und der Ersetzungs-Verbund, nicht die Isolation gegen
  alle übrigen Achsen. Was der Anwender nicht kann: eine Veröffentlichung originalgetreu isoliert durchmessen
  (Abschnitt~\ref{sec:outlook}).
  \item \textbf{Keine automatische Auswahl der besten Konfiguration:} Die heuristik-gesteuerte
  Vollautomatisierung (Klassifikator und Binary-Packager) ist nicht implementiert; vorhanden sind das
  messwert-getriebene Inkrement der Binary-Auswahl (Rangierung je Metrik aus der Auswertungs-Ausgabe)
  und die laufzeit-dynamische Cache-Line-Politik. Was der Anwender nicht kann: aus Lastprofil- und
  Hardware-Merkmalen automatisch eine Achsen-Konfiguration empfehlen lassen; die Empfehlung bleibt
  manuelle Ableitung aus der Auswertung (Abschnitt~\ref{sec:outlook}).
\end{itemize}
\end{nurlang}
\begin{nurkurz}
%% INTERN-272: KA-120
\section{Limitierungen}\label{sec:limitations}
Die folgenden Grenzen beschreiben, was ein Anwender mit dem System nicht erreichen kann; Vorbehalte
einzelner
Messreihen führt der Mess-Anhang (Tabelle~\ref{tab:le:limitierung}), Implementierungsstände
Kapitel~\ref{ch:impl}, den Umsetzungsstand von Bau und Smoke-Reihe Abschnitt~\ref{sec:answers}.
Maschinenspezifische empirische Optimierung
ist der Wesenszug cache-bewusster Verfahren (ATLAS~\cite{whaley2001atlas}, FFTW~\cite{frigo1998fftw};
cache-oblivious Verfahren~\cite{frigo1999cache} als Gegenpol), und ein Apparat, der sie zum Prinzip macht,
erbt ihre Kosten; die Modul-Bodies sind keine Grenze (der Codegenerator des Mess-Pfads emittiert je
Permutation die achtzehn voll qualifizierten Achsen-Typen samt Umbrella-Include,
Abschnitt~\ref{sec:impl-pipeline}).
\begin{itemize}
  \item \textbf{Versionswechsel eines Achsen-Algorithmus:} Jede Tier-Binary, deren Overlay-Schnitt die
  geänderte Datei enthält, wird zu einer neuen Identität (Versions-Lock,
  Abschnitt~\ref{sec:impl-contracts});
  ihr Lager-Bestand ist neu zu erzeugen. Was der Anwender nicht kann: einen einzelnen Algorithmus
  austauschen
  und die Messungen der übrigen Binaries behalten.
  \item \textbf{Lager je Maschinentyp und Hardware-Konfiguration:} Messwerte sind an die Hardware-Identität
  der messenden Maschine geschlüsselt (Bestand~2, Abschnitt~\ref{sec:lager}) --- der Preis der
  maschinenspezifischen Optimierung. Was der Anwender nicht kann: Messwerte einer Maschine auf eine andere
  übertragen.
%% INTERN-272: KA-121
  \item \textbf{Deckung der Hardware-Zähler je Mikroarchitektur:} Erhoben werden vier generische Zähler
  (L1-, Last-Level-, dTLB-, Branch-Misses über den Linux-Systemaufruf \texttt{perf\_event\_open}) und zwei
  modell-gebundene Rohzähler (L2-Demand-Misses, Fills aus einem fremden CCX als Kohärenz-Beobachtbare) nur
  mit gegengeprüfter Kodierung (AMD Zen~5); sonst tragen die Felder eine explizit kodierte Nicht-Erhebung
  (n/a mit Grund, nie eine Null: geplant fehlende Daten verzerren
  nicht~\cite{graham2006planned,graham2009missing}); das Last-Level-Ereignis scheitert auf Zen~5 beim Öffnen
  (keine generische Abbildung), RAPL ist best-effort und ohne Zonen-Leserecht leer, IPC ist als Kategorie
  deklariert, nicht erhoben; in der Smoke-Reihe sind die Hardware-Spalten n/a
  (Tabelle~\ref{tab:le:limitierung}, Zeile~1). Was der Anwender nicht kann: L2- und Kohärenz-Zähler auf
  anderen Modellen als Zen~5 erhalten (Abbildung~\ref{fig:measurement-triangle}).
  \item \textbf{Hardware-Passung ohne automatische Prüfung:} Die Hardware-Erkennung liegt in drei
  Funktionen vor: der Erhebung, der Übersetzungszeit-Factory und dem Laufzeit-Verdikt; Runner-Anbindung,
  Rückschrieb, Provenienz-Roundtrip und Passungs-Stempel fehlen. Was der Anwender nicht kann: die Passung
  zwischen Maschine und Binary
  automatisch prüfen lassen (Abschnitte~\ref{sec:hw-probe} und~\ref{sec:impl-system-axes}).
  \item \textbf{Ausgabe ohne Blatt-Auffächerung:} xlsx ist der Standard, die CSV-Form ihr aus ihr erzeugtes
  Gegenstück (Abschnitt~\ref{sec:eval-identity}); Blatt-Auffächerung und CSV-Baumform sind vorgesehen, nicht
  implementiert. Was der Anwender nicht kann: Ergebnisse je Unter-Achsen-Permutation als eigene Blätter
  erhalten oder die CSV-Ablage als Baum erhalten.
  \item \textbf{Permutationsraum und Budget:} Ein vollständiger Bau des golden-Referenzkatalogs kostet je
  System-Permutation Lager und Zeit in dem Umfang, den der Umsetzungsstand beziffert
  (Abschnitt~\ref{sec:answers}); der Planer rechnet die Mächtigkeit je Profil und bietet im Subkommando
  \texttt{simulate} Budget-Eingänge. Was der Anwender nicht kann: den
  vollen Produktraum aller Achsen-Kataloge materialisieren --- er wählt Kataloge oder Teilmengen, und die
  Wiederholungszahl je Messung bleibt eine Abwägung von Aufwand und Konfidenz~\cite{kalibera2013rigorous}.
  \item \textbf{Kein Laufzeit-Wechsel zwischen Tier-Binaries:} Die Hybrid-Stufe ist als Klassifikation,
  Dock-Schicht, XML-Parser und Binary-Proxy implementiert (Abschnitt~\ref{sec:impl-hybrid-lager}); Export,
  Eviction und Router fehlen. Was der Anwender nicht kann: ein Hybrid-Tier betreiben, das zur Laufzeit
  zwischen vorgeladenen Binaries wechselt; auch die Zustands-Übergabe beim Wechsel --- der Memento-Abgleich
  mit den zwei Release-Artefakten single und hybrid im Lager --- ist entworfen, nicht implementiert.
  \item \textbf{Betriebssystem-Regime:} \texttt{operating\_system} ist dreiwertig deklariert; gemessen wird
  unter root-Linux, das zweite Regime (immutables Betriebssystem) ist nicht umgesetzt
  (Abschnitt~\ref{sec:eval-platforms}). Was der Anwender nicht kann: Messreihen unter Windows oder macOS
  oder
  ohne perf-Rechte (\texttt{perf\_event\_paranoid}, RAPL-Zonen) erheben.
  \item \textbf{Keine isolierte Nachmessung einer einzelnen Veröffentlichung:} Der Betriebsmodus, der genau
  eine Veröffentlichung als geschlossenes Achsen-Set anbietet, ist nicht implementiert. Was der Anwender
  nicht
  kann: eine Veröffentlichung originalgetreu isoliert durchmessen (Abschnitt~\ref{sec:outlook}).
  \item \textbf{Keine automatische Auswahl der besten Konfiguration:} Die Vollautomatisierung ist nicht
  implementiert; vorhanden sind das messwert-getriebene Auswahl-Inkrement und die laufzeit-dynamische
  Cache-Line-Politik. Was der Anwender nicht kann: aus Lastprofil-
  und Hardware-Merkmalen automatisch eine Achsen-Konfiguration empfehlen lassen; die Empfehlung bleibt
  manuelle Ableitung aus der Auswertung (Abschnitt~\ref{sec:outlook}).
\end{itemize}
\end{nurkurz}

\begin{nurlang}
\section{Ausblick}\label{sec:outlook}
Kurzfristig stehen die folgenden Schritte an; implementierte Teile treten mit ihrem offenen Rest auf,
entworfene Teile mit ihrem Umsetzungsstand:
\begin{enumerate}[label=(\arabic*)]
  \item der vollständige Bau nach dem deklarierten Umfang (Abschnitt~\ref{sec:eval-build-scope}) und die
  Abgabe-Messung über den 320er-Katalog;
  \item der End-to-End-Lauf der drei Pflicht-Messreihen auf den registrierten Messmaschinen samt Erzeugung der
  Mess-Diagramme;
  \item der Lager-Vollausbau: die Binaries \emph{mit} und \emph{ohne} Messfühler als zwei eingelagerte
  Release-Artefakte und die synthetisierten Funktions-Kurven als Dateien --- die Messung ist immer
  eingeschaltet, im Lager wird nur Fehlendes nachgemessen (Abschnitt~\ref{sec:lager}); der Ausbau ist
  nicht implementiert;
  \item die Hardware-Erkennung: Anbindung an Runner und Prüf-Dock, Ergebnis-Rückschrieb,
  Provenienz-Roundtrip und Passungs-Stempel;
  \item die Blatt-Auffächerung der xlsx-Mappe je Unter-Achsen-Permutation plus Info-Blatt --- der
  Mappen-Writer ist implementiert, die Auffächerung ist nicht implementiert;
  \item die Gegenprüfung der Intel-Rohzähler auf prod2, damit L2- und Kohärenz-Zähler auch dort erhoben
  werden;
  \item das zweite Betriebssystem-Regime (immutables Betriebssystem) und der Windows-Zählertreiber;
  \item der Paper-Isolations-Modus aus FF2.c: ein Betriebsmodus, der die Entwurfsbestandteile genau
  einer Veröffentlichung als geschlossenes Achsen-Set isoliert anbietet, um sie --- bei freistehenden
  Mess- und System-Achsen --- originalgetreu durchzumessen;
  \item der Profil-Anschluss der Organ-Meta-Meta-Achse \texttt{disk\_io} und der Permutationsbau über
  die Scheduling-Unter-Achse;
  \item die Voll-Permutation der Workloads gegen alle Lebewesen-Binaries (FF2);
  \item die Auswertung des H2-Kriteriums nach dem Lauf der Reihe~A\@.
\end{enumerate}

\paragraph{Heuristik-Automation und Binary-Delivery.} Der in der Aufgabenstellung (Teilaufgabe~7)
geforderte Mechanismus, der aus den Messergebnissen automatisch die beste Konfiguration ableitet, ist
in seiner heuristik-gesteuerten \emph{Voll}automatisierung nicht implementiert und Gegenstand künftiger
Arbeit; ein erstes Inkrement der automatischen Auswahl und des Binary-Versands liegt als
eigenständiges Werkzeug (\texttt{best\_binary\_selector}) vor, das je Metrik die beste Permutation aus
der Auswertungs-Ausgabe rangiert und ausliefert --- die xlsx-Mappe ist die Ausgabe und der Standard, die
CSV-Form ihr aus ihr erzeugter, geordneter Baum aus CSV-Dateien (ein Ordner je Blatt, kein einzelnes
Kind-Dokument), und das Werkzeug liest die CSV-Projektion, die die Naht als eine flache Datei schreibt (die
Baumform ist vorgesehen, aber nicht implementiert). Neben der Auswahl (Zusage~1) umfasst
derselbe Werkzeug-Strang die
laufzeit-dynamische Cache-Line-Anpassung (Zusage~2), die für die Cache-Line-Bewusstheit der Arbeit
einschlägig ist: Aus einem ausgewerteten Profil-Aggregat leitet der Politik-Selektor die Cache-Line-Politik
ab. Der Vollausbau setzt drei Komponenten voraus; die zweite liegt vor, und zur ersten wie zur dritten
stellt das Werkzeug ein erstes, rein messwert-getriebenes Inkrement: (i) einen
\emph{Heuristik-/ML-Klassifikator}, der aus den Lastprofil- und Achsen-Features (Operationsmix,
Schlüssel-Charakteristik, Hardware-Kennung) auf die empfohlene Achsen-Konfiguration abbildet; (ii) den
implementierten \emph{XML-Parser} (Import und Export samt Profil-Lader) für die als
XML-Lastprofil-Ergebnis abgelegten, je Architekturfokus wiederverwendbaren Heuristik-Richtlinien; und
(iii) einen \emph{Binary-Packager}, der den so ausgewählten Gesamtalgorithmus als eigenständige Binary
hinter dem einheitlichen Suchalgorithmus-Interface ausliefert. Voraussetzung für alle drei ist, dass die
End-to-End-Messreihen vollständig vorliegen --- erst die belastbare Datenbasis über alle Permutationen
und Lastprofile macht eine messgetriebene, automatisch gesteuerte Auswahl sinnvoll; ohne diese Automatisierung werden
die Konfigurations-Empfehlungen des Auswertungs-Schemas aus Abschnitt~\ref{sec:eval-results} nach
Vorliegen der Daten manuell abgeleitet.

\medskip
%% INTERN-272: KA-122
Mittel- bis langfristig: die Messung der PIM-Allokator-Familie auf PIM-Hardware\footnote{PIM
(Processing-in-Memory): Speicherbausteine mit eigenen Rechenwerken, etwa UPMEM-DPUs oder Samsungs
HBM-PIM\@; ein PIM-Allokator verlangt solche Hardware.} --- der Allokator A16 ist als Achsen-Baustein
integriert und fällt ohne PIM-Hardware auf die portable ausgerichtete Allokation
zurück~\cite{pimmalloc2026} ---; die formale Verifikation ausgewählter Permutations-Klassen im Stil
von StarMalloc~\cite{reitz2024starmalloc}\footnote{StarMalloc ist in F\textsuperscript{$\ast$} mit
dem Rahmenwerk Steel verifiziert, einer Separationslogik für nebenläufige Programme; die Verifikation
eines Achsen-Bausteins in diesem Stil ist Ausblick.} --- die Allokator-Familie A13 ist integriert, die
Verifikation selbst ist offen ---; die Übertragung des Anatomie-Prinzips auf
Beschleuniger-Speicherhierarchien (GPU), ohne Zusage einer Cluster-Verfügbarkeit; und die
Hybrid-Komponenten Live-Ranking, Zustands-Übergabe und Messfehler-Interpolation
(Abschnitt~\ref{sec:heuristics-modes}). Die Vier-Repository-Architektur (Abschnitt~\ref{sec:repos})
--- Cache-Engine, Prüfling, Super-Repository, Manuskript --- erlaubt, weitere Prüflinge analog zu
PRT-ART einzubringen, ohne die Bausteine-Bibliothek zu ändern: Der Prüflings-Kontrakt der Bau-Konfiguration
prüft den Code-Weg und den Registry-Weg gegeneinander, damit beide Hälften der Integration auf
denselben Ort zeigen. Über die hier behandelten baumartigen Suchstrukturen hinaus ist das
Anatomie-Prinzip über höherwertige Abstraktionen auf die weiteren Strukturklassen aus
Abschnitt~\ref{sec:overview} erweiterbar --- hash-basierte, räumliche und flache Strukturen ---, sodass
derselbe Mess- und Permutationsapparat ohne Neubau auf neue Klassen angewendet werden kann; darin liegt
ein wesentlicher Hebel für künftigen experimentellen Fortschritt.
\end{nurlang}
\begin{nurkurz}
\section{Ausblick}\label{sec:outlook}
Kurzfristig stehen die folgenden Schritte an; implementierte Teile treten mit ihrem offenen Rest auf,
entworfene mit ihrem Umsetzungsstand:
\begin{enumerate}[label=(\arabic*)]
  \item der vollständige Bau nach dem deklarierten Umfang (Abschnitt~\ref{sec:eval-build-scope}) und die
  Abgabe-Messung über den 320er-Katalog;
  \item der End-to-End-Lauf der drei Pflicht-Messreihen samt Diagrammen;
  \item der Lager-Vollausbau: die Binaries mit und ohne Messfühler als zwei Release-Artefakte,
  Kurven-Dateien, nur Fehlendes nachmessen (Abschnitt~\ref{sec:lager});
  \item die Hardware-Erkennung: Runner-Anbindung, Rückschrieb, Provenienz-Roundtrip, Passungs-Stempel;
  \item die Blatt-Auffächerung der xlsx-Mappe je Unter-Achsen-Permutation plus Info-Blatt;
  \item die Gegenprüfung der Intel-Rohzähler auf prod2;
  \item das zweite Betriebssystem-Regime und der Windows-Zählertreiber;
  \item der Paper-Isolations-Modus aus FF2.c;
  \item der Profil-Anschluss von \texttt{disk\_io} und der Permutationsbau über die Scheduling-Unter-Achse;
  \item die Voll-Permutation der Workloads gegen alle Lebewesen-Binaries (FF2);
  \item die Auswertung des H2-Kriteriums nach dem Lauf der Reihe~A\@.
\end{enumerate}

\paragraph{Heuristik-Automation und Binary-Delivery.} Die in Teilaufgabe~7 der Aufgabenstellung geforderte
heuristik-gesteuerte \emph{Voll}automatisierung der Konfigurationswahl ist nicht implementiert; ein erstes
Inkrement liegt als
Werkzeug \texttt{best\_binary\_selector} vor, das je Metrik die beste Permutation aus der
Auswertungs-Ausgabe rangiert und ausliefert (die xlsx-Mappe ist die Ausgabe, die CSV-Form ihr erzeugter
Baum). Neben der Auswahl (Zusage~1) trägt derselbe Strang die laufzeit-dynamische Cache-Line-Anpassung
(Zusage~2). Der Vollausbau braucht (i) einen \emph{Heuristik-/ML-Klassifikator}
von Lastprofil- und Achsen-Features auf die Konfiguration, (ii) den implementierten \emph{XML-Parser} der
Heuristik-Richtlinien und (iii) einen \emph{Binary-Packager} für die gewählte Binary; Voraussetzung sind
vollständige End-to-End-Messreihen, bis dahin werden die Konfigurations-Empfehlungen aus
Abschnitt~\ref{sec:eval-results} manuell abgeleitet.

\medskip
%% INTERN-272: KA-122
Mittel- bis langfristig: die Messung der PIM-Allokator-Familie auf PIM-Hardware (Processing-in-Memory:
Speicherbausteine mit eigenen Rechenwerken; A16 ist integriert~\cite{pimmalloc2026}); die formale
Verifikation ausgewählter Permutations-Klassen im Stil von StarMalloc~\cite{reitz2024starmalloc} (A13 ist
integriert); die Übertragung des Anatomie-Prinzips auf Beschleuniger (GPU); und die Hybrid-Komponenten
Live-Ranking, Zustands-Übergabe und Messfehler-Interpolation (Abschnitt~\ref{sec:heuristics-modes}). Die
Vier-Repository-Architektur (Abschnitt~\ref{sec:repos}) nimmt weitere Prüflinge analog zu PRT-ART auf, ohne
die Bausteine-Bibliothek zu ändern. Über Baumstrukturen hinaus ist das Anatomie-Prinzip
auf die hash-basierten, räumlichen und flachen Klassen aus Abschnitt~\ref{sec:overview} erweiterbar, ohne
den
Apparat neu zu bauen --- ein wesentlicher Hebel für künftigen experimentellen Fortschritt.
\end{nurkurz}

%% INTERN-272: S-103
